\uuid{WIg3}
\exo7id{4175}
\titre{exo7 4175}
\auteur{quercia}
\organisation{exo7}
\datecreate{2010-03-11}
\isIndication{false}
\isCorrection{false}
\chapitre{Fonction de plusieurs variables}
\sousChapitre{Dérivée partielle}
\cpge{mpsi-mp2i.polynomes-fractions-rationnelles.c, pcsi.polynomes.c, ptsi.polynomes.c, mpsi-mp2i.polynomes-fractions-rationnelles.d, pcsi.polynomes.d, ptsi.polynomes.d, bcpst1.polynomes-reels.b, mp-mpi.topologie-espaces-normes.e, pc.espaces-vectoriels-normes.e, psi.espaces-vectoriels-normes.e, pt.fonctions-deux-trois-variables.b}
\module{Analyse}
\niveau{L2}
\difficulte{}

\contenu{
\texte{
Soit $U$ l'ensemble des polynômes à coefficients réels de degré~$n$
et à racines réelles simples.
}
\begin{enumerate}
    \item \question{Montrer que $U$ est ouvert dans~$\R_n[X]$.}
    \item \question{Pour $P\in U$ on note $x_1 < x_2 < \dots < x_n$ les racines de~$P$.
Montrer que l'application $P  \mapsto (x_1,\dots,x_n)$ est de classe $\mathcal{C}^\infty$.}
\end{enumerate}
}
